% Ad Atticum 5.5 --- headnote
% ad-atticum-05-05, ca. 14--15 May 51 BC, written at Venusia

\textit{Cicero to Atticus, written at first light at Venusia
shortly before setting out southward --- the Perseus dateline
gives \textit{ii aut prid. Id. Mai.}, 14 May 51 BC, and the
letter says it was handed off as Cicero left Venusia on the
morning of the Ides itself, 15 May. The opening is unusually
flat: nothing to charge, nothing to tell, no room for joking
--- ``so many cares press on me.'' The note is plainly written
at the rate at which the journey south is being made, day by
day. Atticus's letters are to follow him to Brundisium, where
Cicero will be waiting on his legate Pomptinus through the
date Atticus himself has fixed.}

\textit{The second section signals two pieces of business in the
shorthand the friendship runs in. The first is the
conversations Cicero is about to have at Tarentum with Pompey
on the state of the republic --- Greek \textit{dialogous},
shading the word toward Platonic colour --- which he promises
to write out in full afterwards. The second is the financial
item: an unsettled sum of twenty thousand and eight hundred
thousand sesterces (the meaning of the doubled numeral is the
old crux), which Cicero wants cleared before Atticus leaves
Rome for Epirus. The wider Tullia-dowry and prefects business
of the adjacent letters does not surface here by name; the
\textit{quod auctore te velle coepi adiutore adsequar} at the
close suggests, however, that the financial item is part of
the same Oppius transaction pressed in 5.4.}
